\documentclass[11pt]{article}
\usepackage[T1]{fontenc}
\usepackage{amsmath,amssymb}
\usepackage[margin=1in]{geometry}

\title{Injectivity of the Integer Combination Map}
\author{}
\date{}

\begin{document}
\maketitle

Let $a,b\in\mathbb C$ satisfy
\[
  \Delta=\operatorname{Im}(\overline a b)\ne 0.
\]
We must show that $F:\mathbb Z^2\to\mathbb C$, defined by
$F(m,n)=ma+nb$, is injective. In particular, $a\ne0$, since
$a=0$ would make $\Delta=0$.

\section*{Main line}

Suppose $F(v)=F(w)$. Gap 2 reduces this equality to
$(v_1-w_1)a+(v_2-w_2)b=0$. Multiply by $\overline a$ and take
imaginary parts. Gap 1 gives $(v_2-w_2)\Delta=0$, so $v_2=w_2$.
The remaining equation is $(v_1-w_1)a=0$; because $a\ne0$, also
$v_1=w_1$. Thus $v=w$.

\section*{Gaps to prove independently}

\paragraph{Gap 1.}
For all $a,b\in\mathbb C$ and $m,n\in\mathbb Z$,
\[
  \operatorname{Im}\bigl(\overline a(ma+nb)\bigr)
  =n\operatorname{Im}(\overline a b).
\]
The integer coefficients are viewed as real scalars. This identity
isolates the coefficient of $b$.

\paragraph{Gap 2.}
For all $a,b\in\mathbb C$ and $v,w\in\mathbb Z^2$, if
$v_1a+v_2b=w_1a+w_2b$, then
$(v_1-w_1)a+(v_2-w_2)b=0$.

\section*{Allowed facts}

Use elementary complex arithmetic, that $\overline a a$ is real, and
the zero-product property for real and complex numbers. No result
asserting injectivity or linear independence of $a,b$ may be assumed.

\end{document}
